\subsection{Metodología de la Investigación}



\subsubsection{Marco legal del proyecto}
El proyecto se enmarca en normativa técnica y jurídica aplicable al tratamiento de información, prestación de servicios digitales y seguridad de la información.

\paragraph{Normativa internacional}
\begin{itemize}
    \item \textbf{ISO/IEC 27001}: Sistema de gestión de seguridad de la información (SGSI), aplicable para definir políticas, controles y gestión de riesgos de la plataforma (\cite{iso27001}).
    \item \textbf{ISO/IEC 27017}: Controles de seguridad para servicios en la nube, relevante para la operación del modelo SaaS (\cite{iso27017}).
    \item \textbf{ISO/IEC 27018}: Protección de datos personales en servicios cloud públicos, útil para el tratamiento responsable de datos de usuarios (\cite{iso27018}).
    \item \textbf{ISO/IEC 27701}: Extensión de privacidad sobre ISO/IEC 27001, orientada a la gestión de información personal (PIMS) (\cite{iso27701}).
    \item \textbf{IEEE 29148}: Lineamientos para la ingeniería de requisitos de sistemas y software, aplicable a la definición y trazabilidad de requisitos funcionales y no funcionales (\cite{ieee29148}).
    \item \textbf{IEEE 1012}: Prácticas de verificación y validación de sistemas y software, aplicables a pruebas y control de calidad del proyecto (\cite{ieee1012}).
\end{itemize}

\paragraph{Normativa nacional (Bolivia)}
\begin{itemize}
    \item \textbf{Constitución Política del Estado}: reconoce el derecho a la privacidad, intimidad y protección de datos personales, principios base para la gestión de información de usuarios (\cite{cpe2009}).
    \item \textbf{Ley N\textsuperscript{o} 164 (Ley General de Telecomunicaciones, Tecnologías de Información y Comunicación)}: establece el marco general para servicios TIC y uso de medios electrónicos en el país (\cite{ley164}).
    \item \textbf{Decreto Supremo N\textsuperscript{o} 1793}: reglamenta la Ley N\textsuperscript{o} 164 en aspectos operativos vinculados a telecomunicaciones y TIC (\cite{ds1793}).
    \item \textbf{Ley N\textsuperscript{o} 453 (Ley General de los Derechos de las Usuarias y los Usuarios y de las Consumidoras y los Consumidores)}: aplicable a la calidad de atención, información transparente y mecanismos de reclamo en servicios digitales (\cite{ley453}).
    \item \textbf{Ley N\textsuperscript{o} 1080 (Ciudadanía Digital)}: regula mecanismos de interacción digital con validez en trámites y servicios electrónicos, relevante para autenticación e intercambio digital de información (\cite{ley1080}).
\end{itemize}

Adicionalmente, para la integración con WhatsApp y Telegram, el sistema deberá cumplir los términos de servicio, políticas de plataforma y condiciones de uso de sus APIs oficiales, especialmente en materia de consentimiento, envío de mensajes y tratamiento de datos (\cite{whatsapp_business_terms,telegram_terms}).
